We presented an execution-aware A* search framework for cross-exchange stablecoin arbitrage that integrates real-world trading costs—fees, slippage, withdrawal costs, gas fees, transfer latency, and exchange reliability—directly into the pathfinding process. Three domain-specific guidance heuristics ($h_1$, $h_2$, $h_4$) and a multi-start search-control strategy ($S_1$) were proposed to guide the search toward profitable and executable paths through a graph constructed from 12 major centralized exchanges.

Empirically, we observe profitable path planning with net returns of approximately 0.04\% after all fees under realistic execution assumptions, with a portfolio of at least \$10,000 required to realize meaningful profits. Profitable paths consistently exhibit a three-hop structure (transfer--trade--transfer or trade--transfer--trade), a finding validated by the brute-force 3-hop enumeration baseline, which discovers the same optimal paths but requires evaluating substantially more candidate paths, with the gap widening as graph size increases. Monte Carlo backtesting across 500 independent market snapshots confirms the robustness of these results across varying market conditions and order sizes.

Our results demonstrate that incorporating execution constraints directly into heuristic search bridges the gap between theoretical arbitrage detection and practical trading systems. By modeling real-world frictions within the graph structure itself, we establish execution-aware pathfinding as a viable foundation for scalable, cross-exchange arbitrage under realistic market conditions.

\subsection{Future Work}

Future work will focus on strengthening both the system architecture and the dynamic market model. One key direction is \emph{adaptive API caching}: our results show that short caching windows in the interval $(0, 60]$ seconds significantly reduce traversal latency by limiting blocking data requests, but reduce data freshness. A dynamic caching policy that adjusts refresh intervals per node based on volatility, liquidity depth, or traversal frequency would formalize the trade-off between execution speed and data integrity.

A second extension involves fully decoupling search threads from graph update threads. Although A* instances execute in parallel, API fetch operations can still introduce latency variability. Separating search from data acquisition through non-blocking graph snapshots would allow uninterrupted traversal during high-frequency polling and improve real-time stability.

A third direction is extending the framework to decentralized exchanges (DEXs) and automated market makers (AMMs), where slippage is deterministic via bonding curves and can be precisely modeled within the heuristic. This would enable cross-venue arbitrage spanning both centralized and decentralized markets.

\subsection{Limitations}

The current framework has not executed live trades on the blockchain for the presented data. Although we use real market pricing and fees, the profitability of a planned path may change between planning and execution. In practice, stablecoin markets rebalance under execution pressure, and order-book depth can be consumed by other market participants during the execution window.

Incorporating higher-order response models—such as stochastic signals that approximate the evolution of the order book under incremental trades—would allow the algorithm to anticipate how the execution of one trade alters subsequent exchange rates. This would shift the system from static arbitrage detection to predictive execution-aware planning under dynamic market conditions.

Additionally, our heuristics rely on exchange reliability scores derived from historical uptime and community sentiment, which may not capture sudden exchange failures or regulatory actions. Future work should explore real-time reliability estimation from on-chain and off-chain signals.
